\documentclass[aspectratio=169]{beamer}

% Shared Beamer settings for Elements of AIML lectures (UPES).
% Keep this file free of \documentclass to allow reuse via \input.

% Allow lecture sources to use shared theme/assets without copying them into each lecture folder.
% NOTE: These paths are relative to `Unit-xx/Lecture-yy/latex/` (and `Lab/Experiment-xx/latex/`).
\makeatletter
\def\input@path{{../../../_shared/}{../../../_shared/UPES-Beamer-Theme/}}
\makeatother

\usetheme{UPES}
\setbeamertemplate{navigation symbols}{}

\usepackage{amsmath}
\usepackage{amssymb}
\usepackage{booktabs}
\usepackage{graphicx}
\usepackage{hyperref}
\usepackage{xcolor}
\usepackage{listings}

% Code listings (general-purpose).
\lstset{
  basicstyle=\ttfamily\small,
  keywordstyle=\color{blue},
  commentstyle=\color{gray},
  stringstyle=\color{teal},
  showstringspaces=false,
  columns=fullflexible,
  breaklines=true
}

\usepackage{tikz}
\usetikzlibrary{arrows.meta,positioning,shapes.geometric,calc}

% Per-slide source line (references shown on the slide itself).
\newcommand{\slidesources}[1]{%
  \vfill
  {\tiny\color{gray}\raggedright\sloppy\parbox{\linewidth}{Sources: #1}\par}%
}

% Unit-wise source strings for this subject (used by slides/notes).
% Path is relative to the lecture `latex/` folder (the compilation CWD).
% Source strings for Elements of AIML (used by slides + notes).

\newcommand{\AIMLRepoURL}{https://github.com/tali7c/Elements-of-AIML}

\newcommand{\AIMLLabSources}{%
Provost and Fawcett, \textit{Data Science for Business}.;%
McKinney, \textit{Python for Data Analysis}.;%
WEKA Documentation (Waikato Environment for Knowledge Analysis), accessed 2026-02-28.;%
Matplotlib and Seaborn documentation, accessed 2026-02-28.%
}

\newcommand{\AIMLUnitISources}{%
Rich and Knight, \textit{Artificial Intelligence}, McGraw Hill, 2017.;%
Russell and Norvig, \textit{Artificial Intelligence: A Modern Approach} (4e), Ch. 1--2 (Introduction, Intelligent Agents).;%
Poole and Mackworth, \textit{Artificial Intelligence: Foundations of Computational Agents} (2e), selected sections.%
}

\newcommand{\AIMLUnitIISources}{%
Russell and Norvig, \textit{Artificial Intelligence: A Modern Approach} (4e), Ch. 7--9 (Logical Agents, First-Order Logic, Inference).;%
Huth and Ryan, \textit{Logic in Computer Science} (2e), selected sections on propositional and first-order logic.;%
Genesereth and Nilsson, \textit{Logical Foundations of Artificial Intelligence}, selected chapters.%
}

\newcommand{\AIMLUnitIIISources}{%
Mueller and Massaron, \textit{Machine Learning for Dummies}, For Dummies, 2016.;%
Mitchell, \textit{Machine Learning}, selected chapters on learning basics and evaluation.;%
Scikit-learn documentation, accessed 2026-02-28.%
}

\newcommand{\AIMLUnitIVSources}{%
Mueller and Massaron, \textit{Machine Learning for Dummies}, For Dummies, 2016.;%
James, Witten, Hastie, and Tibshirani, \textit{An Introduction to Statistical Learning} (2e), selected chapters.;%
Scikit-learn documentation, accessed 2026-02-28.%
}

\newcommand{\AIMLUnitVSources}{%
NIST, \textit{AI Risk Management Framework (AI RMF 1.0)}, 2023.;%
Barocas, Hardt, and Narayanan, \textit{Fairness and Machine Learning} (online book), selected sections.;%
Knaflic, \textit{Storytelling with Data}, selected chapters on communicating results.%
}



\title[Elements of AIML]{Elements of AIML}
\subtitle{Unit 02 -- Lecture 03: Clause Form, Resolution \& Unification}
\author{Tofik Ali}
\institute{School of Computer Science, UPES Dehradun}
\date{\today}

\graphicspath{{../images/}{../../../_shared/}{../../../_shared/UPES-Beamer-Theme/}}

\begin{document}

\UPESCoverFrame
\setcounter{framenumber}{0}

\begin{frame}
  \titlepage
  \vspace{0.3em}
  \begin{center}
    \footnotesize Topic focus: \textbf{automated inference} via refutation (resolution)\\[0.25em]
    \footnotesize Repository: \texttt{\AIMLRepoURL}
  \end{center}
  \slidesources{\AIMLUnitIISources}
\end{frame}

\begin{frame}{Quick Links}
  \centering
  \hyperlink{sec:cnf}{\beamerbutton{Clause Form}}\hspace{0.6em}
  \hyperlink{sec:respl}{\beamerbutton{Resolution (PL)}}\hspace{0.6em}
  \hyperlink{sec:resfol}{\beamerbutton{Resolution (FOL)}}\hspace{0.6em}
  \hyperlink{sec:unif}{\beamerbutton{Unification}}\hspace{0.6em}
  \hyperlink{sec:summary}{\beamerbutton{Summary}}
\end{frame}

\begin{frame}{Agenda}
  \tableofcontents
\end{frame}

\section{Clause form (CNF)}
\label{sec:cnf}

\begin{frame}{Learning Outcomes}
  By the end of this lecture, you should be able to:
  \begin{itemize}[<+->]
    \item Convert simple propositional sentences to \textbf{CNF / clause form}.
    \item Explain \textbf{resolution} as a refutation proof method.
    \item Apply resolution to small propositional knowledge bases.
    \item Explain \textbf{unification} and why it is needed for FOL resolution.
  \end{itemize}
\end{frame}

\begin{frame}{Abbreviations \& Symbols}
  \begin{itemize}[<+->]
    \item $KB$: knowledge base (set of sentences/facts/rules)
    \item $\alpha$: query (goal sentence)
    \item $\models$: semantic entailment (true in all models)
    \item $\vdash$: syntactic derivation (provable in a proof system)
    \item $\lnot,\land,\lor$: NOT, AND, OR
    \item CNF: conjunctive normal form; DNF: disjunctive normal form
    \item $\Box$: empty clause (contradiction); $\theta$: substitution; MGU: most general unifier
  \end{itemize}
\end{frame}

\begin{frame}{Refutation View of Entailment}
  Key equivalence (idea used by SAT/resolution):
  \[
    KB \models \alpha \;\;\;\; \Leftrightarrow \;\;\;\; KB \land \lnot \alpha \text{ is unsatisfiable}
  \]
  \vspace{0.4em}
  So to prove $KB \models \alpha$, we:
  \begin{enumerate}
    \item add $\lnot \alpha$ to $KB$,
    \item derive a contradiction (empty clause) using resolution.
  \end{enumerate}
\end{frame}

\begin{frame}{Clause Form (CNF)}
  \begin{itemize}[<+->]
    \item CNF: \textbf{AND} of \textbf{OR} clauses:
    \[
      (a \lor b \lor \lnot c) \land (\lnot a \lor d)
    \]
    \item A \textbf{literal} is a variable or its negation (e.g., $p$ or $\lnot p$).
    \item A \textbf{clause} is a disjunction of literals.
  \end{itemize}
\end{frame}

\begin{frame}{CNF Conversion Recipe (PL)}
  Convert any propositional formula into CNF using these steps:
  \begin{enumerate}[<+->]
    \item Eliminate $\leftrightarrow$ and $\rightarrow$
    \item Push $\lnot$ inward (De Morgan + double negation)
    \item Distribute $\lor$ over $\land$
    \item Split the final $\land$ into a set of clauses
  \end{enumerate}
  \vspace{0.3em}
  \small CNF is a \textbf{representation choice} that makes automated inference systematic.
\end{frame}

\begin{frame}{CNF Conversion (Mini Example)}
  Convert: $(p \rightarrow q) \land p$
  \begin{itemize}[<+->]
    \item Eliminate implication: $(\lnot p \lor q) \land p$
    \item Already in CNF:
    \[
      (\lnot p \lor q) \land (p)
    \]
  \end{itemize}
\end{frame}

\begin{frame}{CNF Conversion (Needs Distribution)}
  Convert: $p \rightarrow (q \land r)$
  \begin{itemize}[<+->]
    \item Eliminate implication:
    \[
      \lnot p \lor (q \land r)
    \]
    \item Distribute $\lor$ over $\land$:
    \[
      (\lnot p \lor q) \land (\lnot p \lor r)
    \]
    \item Clauses: $(\lnot p \lor q)$ and $(\lnot p \lor r)$
  \end{itemize}
\end{frame}

\section{Resolution in propositional logic}
\label{sec:respl}

\begin{frame}{Resolution Rule (PL)}
  From two clauses:
  \[
    (A \lor p),\;\;\; (B \lor \lnot p)
  \]
  we can infer:
  \[
    (A \lor B)
  \]
  \vspace{0.4em}
  Intuition: if either $p$ is true or false, one of the two clauses forces the remaining parts.
\end{frame}

\begin{frame}{Worked Example (PL): $(p\rightarrow q), p \models q$}
  \begin{itemize}[<+->]
    \item $KB = \{p \rightarrow q,\; p\}$, query $\alpha = q$
    \item Add $\lnot q$ and convert to CNF:
    \[
      (\lnot p \lor q),\; (p),\; (\lnot q)
    \]
    \item Resolve $(\lnot p \lor q)$ with $(p)$ gives $(q)$
    \item Resolve $(q)$ with $(\lnot q)$ gives \textbf{empty clause} $\Box$
  \end{itemize}
  Therefore $KB \models q$.
\end{frame}

\begin{frame}{Worked Example (PL): $(p\rightarrow q),(q\rightarrow r),p \models r$}
  \begin{itemize}[<+->]
    \item Add $\lnot r$ and convert to CNF:
    \[
      (\lnot p \lor q),\;(\lnot q \lor r),\;(p),\;(\lnot r)
    \]
    \item Resolve $(\lnot p \lor q)$ with $(p)$ gives $(q)$
    \item Resolve $(\lnot q \lor r)$ with $(q)$ gives $(r)$
    \item Resolve $(r)$ with $(\lnot r)$ gives $\Box$
  \end{itemize}
  This shows \textbf{entailment by contradiction}.
\end{frame}

\section{Resolution in first-order logic}
\label{sec:resfol}

\begin{frame}{FOL Resolution: What's New?}
  \begin{itemize}[<+->]
    \item Variables and quantifiers
    \item Need to convert to a clause form suitable for resolution
    \item When literals ``almost match'', we use \textbf{unification} to make them match
  \end{itemize}
\end{frame}

\begin{frame}{FOL Clause Form Pipeline (High-Level)}
  \begin{enumerate}[<+->]
    \item Remove $\leftrightarrow,\rightarrow$
    \item Push $\lnot$ inward
    \item Standardize variables apart (rename to avoid clashes)
    \item Skolemize $\exists$ (replace with new constants/functions)
    \item Drop $\forall$ (implicit in clauses), distribute, split into clauses
  \end{enumerate}
  \vspace{0.3em}
  \small The output is a \textbf{set of clauses} used by resolution.
\end{frame}

\begin{frame}{Skolemization (Intuition)}
  \begin{itemize}[<+->]
    \item $\exists$ means ``there exists something'' --- introduce a \textbf{new symbol} for it.
    \item Example:
    \[
      \forall x\; \exists y\; Likes(x,y)
    \]
    becomes
    \[
      \forall x\; Likes(x,f(x))
    \]
    \item $f$ is a \textbf{Skolem function} (depends on $x$).
  \end{itemize}
\end{frame}

\begin{frame}{Classic FOL Example: Socrates is Mortal}
  Knowledge:
  \begin{itemize}
    \item $\forall x\; Human(x) \rightarrow Mortal(x)$
    \item $Human(Socrates)$
  \end{itemize}
  Query: $Mortal(Socrates)$
  \vspace{0.4em}
  Refutation: add $\lnot Mortal(Socrates)$ and resolve to contradiction.
\end{frame}

\begin{frame}{Socrates: Clauses + One Resolution Step}
  CNF (clause form):
  \begin{itemize}[<+->]
    \item $\lnot Human(x) \lor Mortal(x)$
    \item $Human(Socrates)$
    \item $\lnot Mortal(Socrates)$ \hfill (negated query)
  \end{itemize}
  \vspace{0.3em}
  Resolve $Human(Socrates)$ with $\lnot Human(x) \lor Mortal(x)$ using $\{x/Socrates\}$ to infer $Mortal(Socrates)$, then derive $\Box$ with $\lnot Mortal(Socrates)$.
\end{frame}

\section{Unification}
\label{sec:unif}

\begin{frame}{What is Unification?}
  Unification finds a substitution that makes two expressions identical.
  \vspace{0.4em}
  \begin{exampleblock}{Example}
    Unify $Likes(x, AI)$ and $Likes(Ali, AI)$
    \[
      \Rightarrow \{x/Ali\}
    \]
  \end{exampleblock}
  We prefer the \textbf{most general unifier (MGU)}.
\end{frame}

\begin{frame}{Occurs-Check (Why some unifications fail)}
  Trying to unify $x$ with $f(x)$ should fail:
  \[
    x = f(x)
  \]
  because it would require an infinite term $x = f(f(f(\dots)))$.
  \vspace{0.4em}
  This check avoids inconsistent substitutions.
\end{frame}

\begin{frame}{Unification Steps (Algorithm Idea)}
  To unify two expressions, repeatedly:
  \begin{enumerate}[<+->]
    \item if they are identical, continue
    \item if one side is a variable, substitute it (if occurs-check passes)
    \item if both are functions/predicates, unify their corresponding arguments
    \item otherwise, fail (no unifier exists)
  \end{enumerate}
  \vspace{0.3em}
  Output is a substitution $\theta$; the best one is the \textbf{MGU}.
\end{frame}

\begin{frame}{Unification Example (MGU)}
  Unify:
  \[
    Knows(x, f(y), y)\;\;\text{and}\;\; Knows(Ali, f(z), z)
  \]
  \begin{itemize}[<+->]
    \item $x \mapsto Ali$
    \item need $f(y)=f(z)$ so $y \mapsto z$
  \end{itemize}
  MGU: $\{x/Ali,\; y/z\}$.
\end{frame}

\section{Summary}
\label{sec:summary}

\begin{frame}{Summary}
  \begin{itemize}
    \item To prove $KB \models \alpha$, check unsatisfiability of $KB \land \lnot \alpha$.
    \item Resolution works on CNF (clauses).
    \item FOL resolution needs unification to match literals with variables.
  \end{itemize}
  \slidesources{\AIMLUnitIISources}
\end{frame}

\begin{frame}{Exit Questions}
  \begin{enumerate}
    \item Why do we convert to CNF before resolution?
    \item State the resolution rule for propositional logic.
    \item Explain unification with an example substitution.
    \item Why does the occurs-check matter?
  \end{enumerate}
\end{frame}

\end{document}
